\documentclass{beamer}
\usetheme{Madrid}
\usecolortheme{default}

\usepackage[utf8]{inputenc}
\usepackage[T1]{fontenc}
%\usepackage[portuguese]{babel}
\usepackage{lmodern}
\usepackage{booktabs}
\usepackage{csquotes}
\usepackage{hyperref}
\usepackage{graphicx}
\usepackage{setspace}
\usepackage{xcolor}
\usepackage{tikz}
\usetikzlibrary{positioning}
\usepackage[style=authoryear-comp]{biblatex}
\addbibresource{biblio.bib}

% -------------------------------------------------------------
% Custom color (Dark Red)
% -------------------------------------------------------------
\definecolor{darkRed}{HTML}{8B0000}

% Adjusting some colors for headings, titles, etc.
\setbeamercolor{structure}{fg=darkRed}
\setbeamercolor{title}{fg=white}     % Title in white
\setbeamercolor{frametitle}{fg=white}
\setbeamercolor{alerted text}{fg=darkRed}
\setbeamercolor{item projected}{fg=white,bg=darkRed}
\setbeamercolor{progress bar}{fg=darkRed,bg=darkRed!50}

\onehalfspacing

% -------------------------------------------------------------
% Insper logo
% -------------------------------------------------------------
\logo{\includegraphics[height=0.8cm]{logo-insper-preto.png}}

% -------------------------------------------------------------
% Title and authors
% -------------------------------------------------------------
\title[Beneath the Surface - LACEA]{\textbf{Beneath the Surface: Investigating Anti-Competitive Behavior in Brazilian Public Procurement}}
\author[Galletta, Genicolo-Martins, Giommoni]{Sergio Galletta, Darcio Genicolo-Martins, Tommaso Giommoni}
\institute[]{ETH Zürich \and Insper \and University of Amsterdam}
\date{\today}

% -------------------------------------------------------------

\begin{document}

%==================================================
\begin{frame}
  \titlepage
\end{frame}

%==================================================
% Slide 1: Motivation
%==================================================

\begin{frame}{Motivation}
  \begin{itemize}
    \item Integrity and efficiency in procurement are vital for the effective use of public funds, economic development, and public trust \cite{BosioEtAl2022}.
    \item Public procurement is a significant component of global GDP (around 12\% overall and up to 20\% in developing countries).
    \item Anti-competitive behavior (collusion, favoritism, bid leakage) increases costs, reduces service quality, and limits innovation.
    \item Detecting anti-competitive behavior is difficult due to its hidden nature and the complexity of procurement processes.
  \end{itemize}
\end{frame}

%==================================================
\begin{frame}{Introduction -- The Paper}
\begin{itemize}
  \item \textbf{Research topic:} Detect signals of anti-competitive behavior in public auctions in S\~{a}o Paulo, Brazil. \pause
  \item \textbf{Data:} 3.8 million procurement auctions from 2010 to 2019, including bid-level and firm-level information. \pause
  \item \textbf{Methods:}
    \begin{itemize}
      \item We use a \textbf{Regression Discontinuity Design (RDD)}, exploiting small differences between narrowly winning and narrowly losing bids (Kawai et al., 2023, \emph{RESTUD}).
      \item Near the cutoff, winning is nearly random; if \textbf{incumbency} or \textbf{backlog} spikes at the threshold, this signals potential anti-competitive behavior.
    \end{itemize} \pause
  \item \textbf{Preview of results:} Evidence consistent with systemic anti-competitive behavior; more work is ongoing on \\ mechanisms and heterogeneity.
\end{itemize}
\end{frame}

%==================================================
\begin{frame}{Our Contribution in the Literature}
\textbf{Relative to Kawai et al. (2023):}
\begin{itemize}
  \item \textbf{Scale and setting:} we take the near-tie RDD approach to a massive, middle-income setting:
  \begin{itemize}
    \item state-wide e-procurement platform (BEC) in S\~ao Paulo,
    \item 3.8 million bid-level records, $\sim$19,000 firms, 1,344 buyers.
  \end{itemize}
  \item \textbf{From incumbency to a richer outcome set:}
  \begin{itemize}
    \item we replicate the \emph{incumbency jump} in near-ties,
    \item and add \textbf{Ratio Won} and \textbf{Backlog 120 days} to study how wins and workload concentrate.
  \end{itemize}
  \item \textbf{Brazilian institutional context:}
  \begin{itemize}
    \item direct evidence from a large emerging economy with documented corruption and enforcement challenges,
    \item both collusive and favoritist mechanisms are a priori plausible.
  \end{itemize}
\end{itemize}
\end{frame}


%==================================================
\begin{frame}{Our Contribution in the Literature}
\textbf{Beyond Kawai and the existing literature:}
\begin{itemize}
   \item \textbf{Ruling out rotation with backlog:}
  \begin{itemize}
    \item classic bid-rotation models predict lower backlog for the designated winner,
    \item we find \emph{higher} backlog for near-tie winners, pointing instead to favoritism/leakage in a repeated e-procurement environment.
  \end{itemize}
  \item \textbf{Turning near-tie RDD into a screening tool:}
  \begin{itemize}
    \item we explicitly frame the design as an \emph{operational risk indicator} for oversight agencies,
    \item mapping where and when the jumps are largest and how they respond to\\ exogenous shocks (random audits, PCC expansion etc.).
  \end{itemize}
\end{itemize}
\end{frame}

%==================================================
\begin{frame}{Anti-Competitive Behavior}
  \begin{itemize}
    \item There are multiple reasons for anti-competitive behavior when the public sector is involved. \vspace{2mm}
    \item \textbf{Collusion:} Agreements among firms that allow them to overcharge for goods and services. \vspace{2mm}
    \item \textbf{Corruption:} Government officials may award contracts to politically connected firms instead of more efficient ones. \vspace{2mm}
  \end{itemize}
\end{frame}

%==================================================
\begin{frame}{Identifying Anti-Competitive Behavior in Pub Proc}
\begin{itemize}
  \item Commonly discussed strategies: \vspace{1mm}
  \begin{enumerate}
    \item \textbf{Incumbency priority}: winners are more likely to be incumbents.
    \item \textbf{Bid rotation}: winners are likely to have lower backlog.
  \end{enumerate} \pause
  \item The \textbf{problem}: there are \textbf{non-collusive, cost-based explanations} for these allocation patterns. \vspace{1mm}
  \begin{itemize}
    \item Bid rotation can arise under competition if \textbf{marginal costs increase with backlog}.
    \item Incumbency advantage can be explained by \textbf{cost asymmetries} among firms or by \textbf{learning-by-doing}.
  \end{itemize}
\end{itemize}
\end{frame}

%==================================================
\begin{frame}{Identifying Anti-Competitive Behavior in Pub Proc}
\begin{itemize}
  \item Kawai et al.'s (2023) approach hinges on the idea that \textbf{closely matched bids imply similar cost structures and winning probabilities}. \vspace{2mm}
  \item Discrepancies in firms' previous outcomes around such near-ties can indicate anti-competitive behavior. \pause \vspace{2mm}
  \item By using a \textbf{Regression Discontinuity Design (RDD)}, one exploits small differences between narrowly winning and narrowly losing bids. \vspace{2mm}
  \item \textbf{Key assumption:} Firms that bid just above or just below the winning margin should be \textbf{very similar in terms of cost structures}. \pause \vspace{2mm}
  \item We are interested in: \textbf{(1) Incumbency status} and \\ \textbf{(2) Backlog} of firms bidding around the cutoff.
\end{itemize}
\end{frame}

%==================================================
\begin{frame}{Regression Discontinuity Design (RDD)}
\[
\beta = \lim_{\epsilon \to 0^+} \mathbb{E}\!\left[x_{i,t} \mid \Delta_{i,t} = \epsilon\right]
   - \lim_{\epsilon \to 0^-} \mathbb{E}\!\left[x_{i,t} \mid \Delta_{i,t} = \epsilon\right]
\]
where:
\begin{itemize}
  \item \( \Delta_{i,t} \): Difference between the bid of firm \( i \) and the most competitive alternative bid at time \( t \).
  \item \( x_{i,t} \): Outcome of interest (incumbency or backlog).
\end{itemize}
\end{frame}

%==================================================
\begin{frame}{Institutional Background}
\begin{itemize}
  \item S\~{a}o Paulo is the wealthiest and most populous state in Brazil (about 46 million inhabitants), contributing over 30\% of national GDP. \vspace{2mm}
  \item Major industries include manufacturing, finance, agriculture, and services. \pause \vspace{2mm}
  \item S\~{a}o Paulo has one of the largest public procurement markets in Brazil, handling billions of reais in contracts annually. \vspace{2mm}
  \item The large scale of procurement creates opportunities for anti-competitive behavior (e.g., a detected school meal cartel was fined over R\$340 million / US\$60 million).
\end{itemize}
\end{frame}

%==================================================
\begin{frame}{Institutional Background}
\begin{itemize}
  \item S\~{a}o Paulo uses the \textbf{Bolsa Eletr\^onica de Compras (BEC)} platform for public procurement. \vspace{2mm}
  \item Two types of tenders: \vspace{2mm}
  \begin{itemize}
    \item \textbf{\textcolor{red}{Sealed-bid tenders (convite)}}: suitable for smaller contracts (under R\$176{,}000, approximately US\$31{,}000). \vspace{2mm}
    \item \textbf{Multi-round descending auctions (preg\~{a}o)}: for larger contracts, including a reverse-auction phase. \vspace{2mm}
  \end{itemize}
\end{itemize}
\end{frame}

%==================================================
% Data
%==================================================

\begin{frame}{Data}
\begin{itemize}
  \item Data from the BEC platform, 2009--2019. \vspace{2mm}
  \item \textbf{3.8 million bid-level records}, covering around 19{,}000 firms and 1{,}344 public buyer units. \vspace{2mm}
  \item Key variables: \vspace{1mm}
  \begin{itemize}
    \item \textbf{Incumbency}: whether a firm won the previous auction it participated. \vspace{1mm}
    \item \textbf{Ratio Won}: share of auctions won in the previous months. \vspace{1mm}
    \item \textbf{Backlog}: total amount of work (value of contracts) a firm has accumulated in the previous months. \vspace{1mm}
    \item \textbf{Margin of Victory}: percentage difference between the lowest and second-lowest bids. \vspace{1mm}
    \item \textbf{Firm characteristics}: age, legal status, and additional \\ controls to be collected.
  \end{itemize}
\end{itemize}
\end{frame}

\begin{frame}{Data}
  \begin{figure}
    \centering
    \includegraphics[width=0.7\textwidth]{total_amount_awarded_per_year (1).pdf}
    %\caption*{Total amount awarded per year (BEC platform, 2009--2019).}
  \end{figure}
\end{frame}

%==================================================
% Graphical Results
%==================================================

\begin{frame}{Graphical Results -- Incumbency Advantage}
  \begin{figure}
    \centering
    \includegraphics[width=0.8\textwidth]{incumbent.png}
    %\caption{Probability of incumbency around the bid margin (near ties).}
  \end{figure}
\end{frame}

\begin{frame}{Graphical Results -- Ratio Won}
  \begin{figure}
    \centering
    \includegraphics[width=0.8\textwidth]{share_won.png}
    %\caption*{Share of auctions won in the last 90 days around the bid margin.}
  \end{figure}
\end{frame}

\begin{frame}{Graphical Results -- Backlog 120 days}
  \begin{figure}
    \centering
    \includegraphics[width=0.8\textwidth]{backlog.png}
    %\caption*{Standardized backlog (120 days) around the bid margin.}
  \end{figure}
\end{frame}


%==================================================
\begin{frame}{Why This Is \textbf{Not} Classic Bid Rotation}
\textbf{Classic rotation cartels (theory predictions):}
\begin{itemize}
    \item Winners tend to have \textbf{lower backlog} (cartel avoids overloading one member).
    \item Recent wins are \textbf{evenly spread} across participants.
    \item Timing does \textbf{not} systematically favor the designated winner.
\end{itemize}

\pause

\textbf{Our near-tie evidence shows the opposite:}
\begin{itemize}
    \item Winners-by-a-hair have \textbf{higher backlog} (120-day workload).
    \item Winners-by-a-hair have \textbf{higher Ratio Won} (win concentration).
    
\end{itemize}

\pause

\begin{itemize}
    \item Conclusion: The joint pattern rules out classic rotation and points instead to \textbf{favoritism, information leakage, or timing coordination}.
\end{itemize}

\end{frame}

%==================================================
% Numerical Estimates
%==================================================

\begin{frame}{Numerical Estimates}
\begin{center}
\begin{tabular}{lccc}
\hline\hline
 & Incumbent & Ratio Won & Backlog 120 \\
 & (1) & (2) & (3) \\
\hline
\\[-4pt]
Winning Auction & 0.033*** & 0.016*** & 0.027*** \\
 & (0.004) & (0.004) & (0.009) \\
\\[-4pt]
Observations & 146{,}999 & 76{,}572 & 89{,}307 \\
Bandwidth & 0.01 & 0.01 & 0.01 \\
Mean DV & 0.672 & 0.494 & 0 \\
SD DV   & 0.469 & 0.326 & 1 \\
\hline\hline
\end{tabular}
\end{center}
\end{frame}

%==================================================
% Balance checks
%==================================================

\begin{frame}{Balance Checks}
\begin{center}
\begin{tabular}{lccc}
\hline\hline
 & Firm Age & Limited liability & Buyer/vendor \\
 &          &                   & same municipality \\
 & (1) & (2) & (3) \\
\hline
\\[-4pt]
Winning Auction   & 0.026 & 0.002 & -0.003 \\
                  & (0.077) & (0.004) & (0.002) \\
\\[-4pt]
Observations & 208{,}732 & 208{,}984 & 208{,}984 \\
Bandwidth    & 0.01 & 0.01 & 0.01 \\
Mean DV      & 10.82 & 0.571 & 0.200 \\
SD DV        & 10.43 & 0.494 & 0.400 \\
\hline\hline
\end{tabular}
\end{center}
\end{frame}

%==================================================
% Future analysis
%==================================================

\begin{frame}{Future Analysis}
\begin{itemize}
  \item Additional balance checks on other firm characteristics (size, experience, participation intensity, geography).
  \item Analysis by type of good, buyer, year, and municipality. \vspace{2mm}
  \item Estimate the $\beta$s and study how they respond to exogenous shocks to the environment:
  \begin{itemize}
    \item Effects of exogenous audits (i.e., randomly assigned audits).
    \item Organized crime effects (e.g., geographical expansion of PCC --- Primeiro Comando da Capital).
  \end{itemize}
\end{itemize}
\end{frame}

%==================================================
% Conclusion
%==================================================

\begin{frame}{Conclusion}
  \begin{itemize}
    \item Narrowly winning firms show higher incumbency and larger backlogs.
    \item Winners and near-losers are similar in age and legal structure; covariate balance is good.
    \item Overall, no evidence of classic bid rotation.
    \item Results are consistent with the presence of systemic anti-competitive behavior (favoritism and/or bid leakage).
  \end{itemize}
\end{frame}

\begin{frame}
  \begin{center}
    \Huge Thank you! \\[8pt]
    \large Questions and comments are very welcome.
  \end{center}
\end{frame}

\end{document}